\chapter{Integraciones}
\label{sec:man-integraciones}

\section{Azure Sentinel}

\subsection{Ingesta streaming}

\emph{Integraciones \textrightarrow{} Sentinel \textrightarrow{}
Conectar}. Te pide:

\begin{itemize}
    \item \code{AZURE_TENANT_ID}, \code{AZURE_CLIENT_ID},
    \code{AZURE_CLIENT_SECRET} — para un Service Principal con
    permisos \emph{Reader} en el workspace.
    \item \code{SENTINEL_WORKSPACE_ID}.
    \item Una query KQL (ej.\ \code{SigninLogs | where
    TimeGenerated > ago(24h)}).
\end{itemize}

Click \emph{Conectar} y la sesión empieza a recibir eventos en
vivo. Podés cerrar la conexión con \emph{Desconectar}.

\subsection{Publish IOCs}

Requiere \code{GRAPHHUNTER_IOC_PUBLISH_ENABLED=1} en el entorno
del daemon. Los IOCs se publican como \emph{Threat Intelligence
Indicators} en el workspace.

\subsection{Estado}

\emph{Integraciones \textrightarrow{} Sentinel \textrightarrow{}
Status} muestra: conexión activa, eventos recibidos, última query,
errores de API.

\section{MCP (Model Context Protocol)}

GraphHunter expone su API como tools MCP. Esto permite que Cursor,
Claude Code, Claude Desktop u otros clientes compatibles accedan
al grafo vivo desde su asistente.

\subsection{Setup}

En tu cliente MCP (ej.\ \filepath{.cursor/mcp.json}):

\begin{lstlisting}[basicstyle=\ttfamily\small]
{
  "mcpServers": {
    "graphhunter": {
      "command": "node",
      "args": ["/path/to/graph-hunter-mcp/dist/src/index.js"],
      "env": {
        "GRAPHHUNTER_API_URL": "http://127.0.0.1:37891",
        "GRAPHHUNTER_API_TOKEN": "<uuid del app startup>"
      }
    }
  }
}
\end{lstlisting}

Copiá el UUID que el app imprime al arrancar. Sin token, el MCP se
niega a arrancar (salvo que fijes
\code{GRAPHHUNTER_MCP_ALLOW_NO_TOKEN=1}, sólo para dev).

\subsection{Qué podés pedir}

Ejemplos de prompts para el asistente:

\begin{itemize}
    \item ``Buscá todas las IPs externas que se conectaron en la
    última hora.''
    \item ``¿Qué relaciones existen en el grafo entre \code{alice}
    y \code{web-01}?''
    \item ``Corre el hunt \code{User -[Auth]-> Host -[Connect]->
    Host \{k=3\}} y explicame el top 5.''
\end{itemize}

El asistente traduce tu pedido a las tools MCP apropiadas
(\code{search_entities}, \code{run_hunt}, etc.).

\subsection{Debug}

Si el asistente no ve las tools:

\begin{itemize}
    \item Verificar que el MCP server arrancó sin errores
    (\code{GRAPHHUNTER_MCP_DEBUG=1} para logs verbosos).
    \item Verificar que el app está corriendo y el token está
    sincronizado.
    \item Reiniciar el cliente MCP (a veces hay que bajar y subir
    Cursor).
\end{itemize}

\section{AI backends (LLM)}

GraphHunter usa LLM para el carril agéntico (parser\_generate,
revisión de mappings). Backend configurable:

\begin{itemize}
    \item \textbf{Local GGUF}: setear
    \code{GRAPHHUNTER_LOCAL_MODEL_PATH=/ruta/modelo.gguf}. Óptimo
    para ambientes air-gapped o con datos sensibles.
    \item \textbf{Remoto (OpenAI, Anthropic, Google)}: configurar
    las credenciales del proveedor (\code{OPENAI_API_KEY}, etc.) y
    elegir el backend en \emph{Settings \textrightarrow{} LLM
    Backend}.
\end{itemize}

El LLM nunca ve eventos de producción — sólo muestras explícitas
que vos le entregás en el panel agéntico.

\section{GeoIP (MaxMind)}

Para \emph{enrich\_ip}: bajar GeoLite2 o GeoIP2 City de MaxMind y
apuntar \code{GRAPHHUNTER_GEOIP_DB=/ruta/GeoLite2-City.mmdb}. El
enriquecimiento es opt-in por llamada; no se auto-enriquece la
ingesta.

\section{Tickets (Jira / ServiceNow)}

\subsection{Jira}

\begin{lstlisting}[basicstyle=\ttfamily\small]
JIRA_BASE_URL=https://tuorg.atlassian.net
JIRA_EMAIL=tu@email.com
JIRA_API_TOKEN=<token>
JIRA_PROJECT_KEY=SEC
JIRA_ISSUE_TYPE=Task        # opcional, default "Task"
GRAPHHUNTER_TICKET_PUBLISH_ENABLED=1
\end{lstlisting}

\subsection{ServiceNow}

\begin{lstlisting}[basicstyle=\ttfamily\small]
SERVICENOW_INSTANCE=tuorg.service-now.com
SERVICENOW_USERNAME=user
SERVICENOW_PASSWORD=<password>
SERVICENOW_TABLE=incident   # opcional, default "incident"
GRAPHHUNTER_TICKET_PUBLISH_ENABLED=1
\end{lstlisting}

Cada llamada a \emph{Crear ticket} requiere \code{confirm:true} o
\code{dry_run:true}. Doble barrera: env var + confirmación.
